\documentclass[11pt,a4paper]{article}
\usepackage[T1]{fontenc}\usepackage[utf8]{inputenc}\usepackage{lmodern}
\usepackage[a4paper,margin=2.2cm]{geometry}
\usepackage{amsmath,amssymb}\usepackage{enumitem}\usepackage{booktabs,longtable,array}
\usepackage[hidelinks]{hyperref}
\setlength{\parindent}{0pt}\setlength{\parskip}{0.4em}
\title{\textbf{OP-CONST T1: the $U(1)$-emergence normalisation and the anchoring
of $\alpha$}\\\large Recon v1 --- proposal-only; no preprint edits}
\author{Per Step 3 review directive: does $q^2/Z_A$ collapse to parameter-only?}
\date{12 June 2026 --- rev.~1.1 (review softenings applied: candidate lemma, conditional channels, explicit C1--C4)}
\begin{document}\maketitle

\section*{0. Scope and firewall}
T1 asks one question at recon level: in the recorded $U(1)$-emergence construction,
does the physical combination $\alpha_{\rm fs}=q^2/(4\pi Z_A)$ retain
stiffness/background dependence, or does it collapse to a parameter-only quantity?
Sources: \S sec:u1\_emergence (Layers 1--3), \S sec:alpha\_fs, Appendix A.3, and the
common-origin matching relations eq:GN\_S0\_common. No D1 block; no preprint edits;
all statements below are classifications of the recorded structure plus explicitly
labelled derivation chains.

\section{What the record fixes (and what it leaves free)}
\textbf{(a) Canonical Maxwell frame; all freedom in $e$.} Proposition~3 records the
leading gauge kinetic term in canonical normalisation, $-\tfrac14 F_{AB}F^{AB}$
(eq:Maxwell\_action), followed by the explicit statement: ``The gauge coupling $e$
is not fixed by the above argument. It remains a free EFT matching parameter''
(L$\sim$6377--6379). The $Z_A$-frame of \S sec:alpha\_fs is the equivalent
presentation $e=q/\sqrt{Z_A}$: the entire normalisation question lives in one
matching coefficient, exactly as open target (vii) states.

\textbf{(b) Topology quantises the product, not the magnitude.} The compact
holonomy theorem records $\oint\mathcal D_A\theta\,dX^A=2\pi n$ and hence
$e\oint A_A\,dX^A=2\pi n$ (eq:holonomy): the flux lattice quantises the
\emph{product} of coupling and flux. Matter charges are integer multiples on this
lattice (ratios rigid; consistent with the $\mathbb Z_6$ interlock layer), but the
magnitude of $e$ itself is not topologically fixed. \textbf{Consequence: the
numerator route is closed.} $q$ cannot acquire compensating $u$-dependence ---
it is lattice-integer --- so $q^2/Z_A$ cannot collapse to parameter-only via $q$;
any protection must come from the composition of $Z_A$ itself.

\section{The $S_0$ lever (central T1 finding)}
\textbf{(a) The recorded identity.} The common-origin matching relations record
(eq:GN\_S0\_common):
\[
G_N=\frac{c_*^4}{8\pi M_G^2},
\qquad
S_0^{\rm EFT}=\frac{c_\perp\,K_\theta}{2\,m_\varphi^2},
\]
with the simplest quartic realisation stated explicitly: $K_\theta\sim\phi_0^2$,
$m_\varphi=\sqrt{2\lambda}\,\phi_0$, $\phi_0=\sqrt{\mu^2/\lambda}$. Evaluating the
identity in that realisation gives
\[
S_0^{\rm EFT}=\frac{c_\perp}{4\lambda}\quad\text{(parameter-only)},
\]
manifestly consistent with the recorded \emph{structural universality} of $S_0$
across the coherent sector (Level A/B): the $\phi_0^2$ factors cancel.

\textbf{(b) The lever.} Read in reverse, the identity ties the anchoring classes
together: since $S_0$ is recorded as a protected structural invariant, the
identity \emph{would} enforce
\[
\zeta_{K_\theta}\;=\;\zeta_{m_\varphi^2}
\qquad\text{(equal background-tracking exponents)},
\]
up to the kinematic factor $c_\perp$, conditionally on the assumptions of the
candidate lemma below. \textbf{Candidate no-go lemma (asymmetric tracking; not yet
a theorem):} asymmetric tracking of $(K_\theta,m_\varphi^2)$ is incompatible with
the recorded $S_0$-universality \emph{unless another matching factor compensates
it}. The statement is conditional on: $S_0$ remaining a structural invariant on
the cosmological branch; $K_\theta$ and $m_\varphi^2$ being evaluated in the same
background/running context; and the absence of compensating factors in $c_\perp$,
the normalisation, or the matching layer.

\textbf{(c) The reduction.} The $Z_B$-channel question $\zeta_{K_\theta}=?$ thus
collapses to a single object: the anchoring of $m_\varphi^2$ (equivalently, of the
scale at which the simplest-realisation formulas are evaluated). The record writes
$m_\varphi=\sqrt{2\lambda}\,\phi_0$ at the bare VEV --- the P-anchored reading ---
but the very same passage flags that ``the relation between the gradient condensate
scale $u_0$ and the bare potential data remains open (OP-grad)''
(L$\sim$4554--4557). \emph{The decisive $\alpha$-question of OP-CONST therefore
merges into OP-grad}: protection holds iff the evaluation scale of the
stiffness/mass pair is anchored to the bare data rather than to the drifting
cosmological amplitude $u(z)$ --- precisely the $u_0\!\leftrightarrow\!(\mu^2,
\lambda)$ closure that the OP3.1 programme already targets.

\section{The $Z_W$ channel and the $\theta_W$ corollary}
The orientation-sector coefficient is recorded as the induced chain
$\kappa_n=\mathcal C_n u_0^2$ (OP-C$_n$); its anchoring class is the anchoring of
$u_0$ --- the same OP-grad object as in \S2(c). Both electroweak kinetic channels
\emph{may} therefore reduce to the same anchoring scale at leading structural
level, but this requires the explicit inheritance maps for $Z_B$ and $Z_W$: the
two channels could differ in inheritance factors, threshold corrections,
$u_0$-dependence, or normalisation conventions. Correspondingly, the Step-3 ratio
statement is conditional: $\sin^2\theta_W=Z_W/(Z_B+Z_W)$ is protected at leading
order \emph{if} $Z_B$ and $Z_W$ share the same background-tracking factor.

\section{T1 verdict}
\begin{enumerate}[leftmargin=1.9em,itemsep=2pt,label=\textbf{V\arabic*.}]
\item $q^2/Z_A$ does \emph{not} collapse to parameter-only via the numerator: $q$
is lattice-integer and $e$ is recorded as a free matching parameter; the holonomy
quantises $e\times$flux, not $e$. The cancellation route (P1) in its naive form is
closed.
\item Protection of $\alpha$ is nevertheless structurally \emph{available} through
the $S_0$ lever: the recorded identity plus $S_0$-universality \emph{would}
enforce equal tracking of $(K_\theta,m_\varphi^2)$ (candidate lemma, \S2(b)) and
reduce the question to the anchoring of
the evaluation scale, with the parameter-level reading
($m_\varphi=\sqrt{2\lambda}\,\phi_0$) explicitly recorded in the simplest
realisation.
\item Outcome A is a conditional protection-theorem sketch, possible only if all
four of the following hold: (C1)~the inheritance
$Z_B^{(0)}\!\leftarrow\!K_\theta^{\rm eff}$ leaves no uncancelled
$u(z)$-dependence in physical $\alpha$ (the hard core of open target (vii));
(C2)~the $K_\theta/m_\varphi^2$ cancellation is valid on the cosmological branch,
not only in the parameter-level simplest realisation; (C3)~the
$u_0\!\leftrightarrow\!(\mu^2,\lambda)$ closure is P-anchored or sufficiently
protected (OP-grad); (C4)~threshold/running endpoints are P-anchored or cancel.
\emph{Not established, but well-posed.} Outcome C remains the quantified fallback
of Step 3 (channel-dependent $\varepsilon$ bounds).
\item Cross-sector consequence: OP-CONST attaches quantitative acceptance
thresholds to OP-grad and to target (vii): the anchoring of $u_0$ must hold to
$|\partial\ln u_{0,{\rm eff}}/\partial\ln u_{\rm cosm}|\lesssim10^{-5}$ (via
$\zeta_A$) and the $\mathcal I_{\rm EW}$ mechanism to $\lesssim10^{-6}$ (via
$\mu$), on the retained band. The deepest open problem (OP-grad) and the
varying-constants stress test are now one connected programme.
\end{enumerate}

\section{Next targets (proposal-layer, review-gated)}
(T1a)~Compute the $Z_B^{(0)}\!\leftarrow\!K_\theta^{\rm eff}$ inheritance factor
explicitly --- the one remaining derivation of the $\alpha$-chain. (T1b)~On the
OP-grad side, classify the anchoring of $u_0$ within the OP3.1 closure-map work
(direct cross-link to the constraint-relaxation bridge programme).
(T1c)~Attempt to promote the candidate no-go lemma of \S2(b) to a theorem by
closing its three conditional assumptions. (T1d)~The $\mu$-side selection
principle of Step 3B
stands unchanged. No D1/preprint contact before review of this document.
\end{document}
